%========================================================================
\section{Le problème : une QA non reproductible}
% ⏱ ~4-5 min  |  Cœur du CADRAGE ingénieur. Le jury académique note ça.
%========================================================================

\begin{frame}{La QA dans le Delivery repose sur une seule condition}
  % 🎤 ~1 min. Les 3 missions QA, et la condition commune : la reproductibilité.
  \begin{columns}[T]
    \column{.6\textwidth}
    \begin{itemize}
      \item \textbf{Homologation} --- valider une nouvelle feature avant MEP
      \item \textbf{TNR} --- vérifier que l'existant n'est pas dégradé
      \item \textbf{QA au sens large} --- un niveau constant, des process reproductibles
    \end{itemize}
    \column{.4\textwidth}
    \begin{block}{Condition unique}
      Rejouer \alert{les mêmes vérifications}, sur \alert{les mêmes données}, dans \alert{les mêmes parcours}.
    \end{block}
  \end{columns}
\end{frame}

\begin{frame}{Trois défaillances empêchaient cette reproductibilité}
  % 🎤 ~2 min. Le cœur du diagnostic. Une défaillance = une cause racine, pas un symptôme.
  \begin{enumerate}
    \item \textbf{Tickets sous-spécifiés} --- critères d'acceptance flous, maquettes absentes
    \item \textbf{Documentation absente ou obsolète} --- seule la Home documentée, et dégradée
    \item \textbf{Pas de Golden Data} --- pré-prod partagée et modifiée quotidiennement $\Rightarrow$
          le référentiel testé change entre deux passes
  \end{enumerate}
  \vskip 6pt
  % 🎤 « Résultat : tests manuels exhaustifs avant chaque MEP, couverture inégale, régressions imprévisibles. »
\end{frame}

\begin{frame}{Première réponse, et pourquoi elle ne suffit pas}
  % 🎤 ~1 min 30. Montrer une réponse PRAGMATIQUE (QA Checks) + sa limite + la REFORMULATION (= maturité ingénieur).
  \begin{columns}[T]
    \column{.5\textwidth}
    \textbf{Les « QA Checks »}
    \begin{itemize}\footnotesize
      \item Section ajoutée au ticket avant dev
      \item Palliatif aux critères d'acceptance
      \item $\Rightarrow$ moins d'allers-retours dev/QA/PO
    \end{itemize}
    \column{.5\textwidth}
    \textbf{Mais limité par construction}
    \begin{itemize}\footnotesize
      \item Vit dans le périmètre d'un ticket
      \item Aucune vision globale du site
      \item Non automatisable
    \end{itemize}
  \end{columns}
  \vskip 8pt
  \begin{alertblock}{Reformulation (avec mon tuteur) $\to$ 2 décisions de conception indépendantes}
    Quelle \alert{stratégie de documentation} des TNR ? \quad Quel \alert{paradigme d'exécution} ?
  \end{alertblock}
\end{frame}
